% Source: https://artofproblemsolving.com/community/c3h2354668_intresting_proof_i_decided_to_make
% Stable identity: aops-topic-2354668-q01
\textbf{(a)} Induction on $k$. For $k=0$ the only number is $1=2^0$. Suppose the statement holds for $k$, and let $1\le m\le2^{k+2}-1$.
\begin{itemize}
\item The sum of all of $2^0,\dots,2^k$ is $2^{k+1}-1<2^{k+1}$. So a representation of $m$ uses $2^{k+1}$ if and only if $m\ge2^{k+1}$.
\item If $m<2^{k+1}$, its representations are exactly those from $2^0,\dots,2^k$, and by the induction hypothesis there is exactly one.
\item If $m\ge2^{k+1}$, then $r=m-2^{k+1}$ satisfies $0\le r\le2^{k+1}-1$, and the representations of $m$ are $2^{k+1}$ plus a representation of $r$ from $2^0,\dots,2^k$. There is exactly one: the empty sum if $r=0$, and the one given by the induction hypothesis otherwise.
\end{itemize}
This completes the induction. (These representations are the binary expansions. Note that $2^0,\dots,2^k$ are $k+1$ powers, so the binary expansions here have $k+1$ digits, not $k$.) $\blacksquare$

\textbf{(b)} Each term of the expansion comes from choosing, in every factor, either the power of $x$ or the constant: $x^1$ or $1$, $x^2$ or $3$, $x^4$ or $5$. The exponent of the term is the sum of the chosen exponents, a sum of distinct numbers from $2^0,2^1,2^2$. By (a) with $k=2$, each exponent from $1$ to $7$ arises from exactly one choice, and the exponent $0$ from choosing all constants. So there is exactly one term for each power $x^0,\dots,x^7$, and no two terms need to be combined.

For $x^5$: $5=4+1$, so choose $x$, then $3$, then $x^4$, giving the coefficient $1\cdot3\cdot1=\boxed{3}$. Indeed
\[(x+1)\left(x^2+3\right)\left(x^4+5\right)=x^7+x^6+3x^5+3x^4+5x^3+5x^2+15x+15.\]
